\documentclass[11pt,a4paper]{article}
\usepackage[margin=22mm]{geometry}
\usepackage{kotex}
\usepackage{amsmath,amssymb,mathtools,bm}
\usepackage{booktabs,longtable,array}
\usepackage{enumitem}
\usepackage{hyperref}
\usepackage{xcolor}
\usepackage{fancyhdr}
\usepackage{lastpage}
\usepackage{setspace}

\setstretch{1.13}
\setlength{\parskip}{0.35em}
\setlength{\parindent}{0pt}
\setlength{\headheight}{14pt}

\hypersetup{
  colorlinks=true,
  linkcolor=blue,
  urlcolor=blue,
  citecolor=blue,
  pdftitle={Graphite ICA DVA Chapter 1 Development Draft},
  pdfauthor={Codex}
}

\pagestyle{fancy}
\fancyhf{}
\lhead{Chapter 1: Charge-Balance Graphite ICA/DVA}
\rhead{\thepage/\pageref{LastPage}}

\newcommand{\dd}{\mathrm{d}}
\newcommand{\ee}{\mathrm{e}}
\newcommand{\OCV}{\mathrm{OCV}}
\newcommand{\app}{\mathrm{app}}
\newcommand{\drive}{\mathrm{drive}}
\newcommand{\bg}{\mathrm{bg}}
\newcommand{\cell}{\mathrm{cell}}
\newcommand{\ext}{\mathrm{ext}}
\newcommand{\tot}{\mathrm{tot}}
\newcommand{\eff}{\mathrm{eff}}
\newcommand{\eq}{\mathrm{eq}}
\newcommand{\refe}{\mathrm{ref}}
\newcommand{\Ah}{\mathrm{Ah}}
\newcommand{\Cunit}{\mathrm{C}}
\newcommand{\sgn}{\operatorname{sgn}}

\begin{document}

\begin{center}
{\LARGE \textbf{Chapter 1}}\\[0.3em]
{\LARGE \textbf{Thermodynamic Charge Balance and Graphite ICA/DVA}}\\[0.8em]
{\large Development draft}
\end{center}

\tableofcontents
\newpage

\section{Purpose and Scope}

This chapter defines the thermodynamic charge-balance basis for graphite-anode
ICA/DVA modeling.  The internal graphite potential is not prescribed as an
external function \(V_{n,\OCV}(q,T)\).  It is solved from charge conservation
after the external charge coordinate, temperature, and internal transition
states have been specified.

The central chain is
\begin{equation}
Q_{\ext}=Q_{\cell}q
\quad\Longrightarrow\quad
Q_{\cell}q =
Q_{\bg}(V_n,T)+\sum_{j=1}^{N_p}Q_{j,\tot}\xi_j
\quad\Longrightarrow\quad
V_n
\quad\Longrightarrow\quad
V_{n,\app}
\quad\Longrightarrow\quad
\frac{\dd Q_{\ext}}{\dd V_{n,\app}},\;
\frac{\dd V_{n,\app}}{\dd Q_{\ext}} .
\label{eq:chapter1_spine}
\end{equation}

Heat generation, detailed electrochemical interface kinetics, global fitting,
and hysteresis are not solved in this chapter.  They are downstream layers that
must consume the quantities defined here without redefining charge balance.

\section{Measured Charge Coordinate}

Let \(Q_{\cell}\) be the reference cell capacity in coulomb.  If a capacity is
given in ampere-hour, the conversion is
\begin{equation}
Q_{\cell}[\Cunit]=3600\,Q_{\cell}[\Ah].
\end{equation}

The external charge passed through the circuit is
\begin{equation}
Q_{\ext}(t)=\int_{t_0}^{t}|I(t')|\,\dd t',
\qquad
q(t)=\frac{Q_{\ext}(t)}{Q_{\cell}} .
\label{eq:q_definition}
\end{equation}

Under a nonzero constant-current segment,
\begin{equation}
\frac{\dd q}{\dd t}=\frac{|I|}{Q_{\cell}},
\qquad
\frac{\dd}{\dd q}=\frac{Q_{\cell}}{|I|}\frac{\dd}{\dd t}.
\label{eq:q_time_conversion}
\end{equation}
The \(q\)-domain form is not used during rest.  At rest, \(q\) is fixed and the
time-domain relaxation equations must be used.

\section{Transition States and Equilibrium Occupancy}

The graphite staging transitions are represented by progress variables
\(\xi_j\), where \(j=1,\ldots,N_p\).  In the discharge-oriented convention used
here, \(\xi_j\) increases as the associated transition contributes positive
external capacity.

For a dummy voltage argument \(V\), the equilibrium progress is written as
\begin{equation}
\xi_{j,\eq}(V,T)
=
\left[
1+\exp\!\left(
-s_{\xi,j}\frac{V-U_j(T)}{w_j(T)}
\right)
\right]^{-1},
\label{eq:xi_eq_dummy_voltage}
\end{equation}
with default \(s_{\xi,j}=+1\) unless a transition requires the opposite
orientation.  The width \(w_j\) is reserved for this equilibrium voltage width.
If a dimensionless capacity fraction is needed, use
\begin{equation}
a_j=\frac{Q_{j,\tot}}{Q_{\cell}},
\end{equation}
rather than reusing \(w_j\).

\section{Charge Balance and Internal Potential}

Define the charge-balance residual
\begin{equation}
\mathcal G(V;q,T,\bm{\xi};\theta)
=
Q_{\bg}(V,T;\theta)
+\sum_{j=1}^{N_p}Q_{j,\tot}\xi_j
-Q_{\cell}q .
\label{eq:charge_balance_residual}
\end{equation}
The internal graphite potential \(V_n\) is the admissible root
\begin{equation}
V_n=\mathcal V(q,T,\bm{\xi};\theta),
\qquad
\mathcal G(V_n;q,T,\bm{\xi};\theta)=0.
\label{eq:V_root_operator}
\end{equation}

The root must lie in the voltage interval used to define the graphite
thermodynamic model.  A sufficient local conditioning requirement is
\begin{equation}
\frac{\partial Q_{\bg}}{\partial V}(V,T)\ge \epsilon_Q>0.
\label{eq:bg_slope_floor_draft}
\end{equation}
This prevents the charge-balance inversion from becoming singular.

\section{Equilibrium OCV as a Derived Relation}

The equilibrium OCV is not an independent input curve.  It is the special
charge-balance solution obtained when each transition is at equilibrium:
\begin{equation}
Q_{\cell}q
=
Q_{\bg}(V,T)
+\sum_{j=1}^{N_p}Q_{j,\tot}\xi_{j,\eq}(V,T).
\label{eq:ocv_implicit_draft}
\end{equation}
The equilibrium anode OCV is therefore
\begin{equation}
V_{n,\OCV}(q,T)=V
\quad\text{such that Eq.~\eqref{eq:ocv_implicit_draft} holds.}
\label{eq:ocv_root_definition}
\end{equation}

The total capacity normalization is imposed by the same storage relation at the
chosen endpoints.  This couples \(Q_{\bg}\), \(Q_{j,\tot}\), \(U_j\), and
\(w_j\); these parameters should not be fitted as unrelated shape knobs.

\section{Apparent and Driving Voltages}

The internal potential, apparent voltage, and kinetic driving voltage are
distinct:
\begin{align}
V_n &=
\mathcal V(q,T,\bm{\xi};\theta),\\
V_{n,\app} &=
V_n+s_I |I|R_n(q,T,|I|),\label{eq:Vapp_draft}\\
\mathcal A_j &=
s_{\phi,j}F\bigl(V_{n,\drive}-U_j(T)\bigr).
\label{eq:affinity_draft}
\end{align}
The approximation \(V_{n,\drive}\approx V_{n,\app}\) may be used only as a
declared modeling choice.  It should not be combined with an unconstrained
fit of both \(R_n\) and the kinetic prefactors, because both can absorb dynamic
voltage shifts.

\section{Self-Consistent Dynamics}

After \(V_n\) has been defined by Eq.~\eqref{eq:V_root_operator}, the transition
dynamics are
\begin{equation}
\frac{\dd \xi_j}{\dd t}
=
k_j(V_n,q,T,I;\theta)
\bigl[\xi_{j,\eq}(V_n,T)-\xi_j\bigr].
\label{eq:dxidt_draft}
\end{equation}

During a nonzero-current segment this becomes
\begin{equation}
\frac{\dd \xi_j}{\dd q}
=
\frac{Q_{\cell}}{|I|}
k_j(V_n,q,T,I;\theta)
\bigl[\xi_{j,\eq}(V_n,T)-\xi_j\bigr].
\label{eq:dxidq_draft}
\end{equation}

Equivalently, with \(V_n(s)=\mathcal V(s,T(s),\bm{\xi}(s);\theta)\),
\begin{equation}
\xi_j(q)
=
\xi_j(q_0)
+
\int_{q_0}^{q}
\frac{Q_{\cell}}{|I(s)|}
k_j(V_n(s),s,T(s),I(s);\theta)
\bigl[\xi_{j,\eq}(V_n(s),T(s))-\xi_j(s)\bigr]
\,\dd s .
\label{eq:volterra_draft}
\end{equation}
This is the explicit self-consistent problem: the state \(\bm{\xi}\) determines
\(V_n\) through charge balance, and the solved \(V_n\) determines the state
update.

\section{Reference-Closure Strategy}

Equations~\eqref{eq:V_root_operator} and \eqref{eq:volterra_draft} define the
direct coupled problem.  The safest numerical baseline is to solve the
charge-balance root at each step and integrate the state equation directly.

For analytic or fast fitting approximations, the closure pattern used in
Refs.~\cite{Lee2011Propagator,Son2013Rates} can be adapted mathematically.
The transferable idea is not the geminate-pair physics, but the way a
self-consistent integral equation is rewritten so that the unknown feedback
appears as a ratio or correction relative to a simpler reference solution.

Choose a reference path \(\bm{\xi}^{\refe}(q)\) and
\(V_n^{\refe}(q)=\mathcal V(q,T,\bm{\xi}^{\refe};\theta)\).  A natural
thermodynamic reference is the equilibrium path
\begin{equation}
\xi_j^{\refe}(q)
=
\xi_{j,\eq}\!\left(V_{n,\OCV}(q,T),T\right).
\label{eq:eq_reference_path}
\end{equation}
Alternative references, such as a low-rate direct solve or a frozen-feedback
path, may be used if they are stated explicitly.

The dynamic integral can then be written schematically as
\begin{equation}
\xi_j(q)=\xi_j(q_0)+
\int_{q_0}^{q}
F_j^{\refe}(s)\,
\mathcal C_j(s;\bm{\xi},\bm{\xi}^{\refe})\,\dd s,
\label{eq:closure_schematic}
\end{equation}
where \(F_j^{\refe}\) is the reference integrand and
\(\mathcal C_j\) is the dimensionless feedback correction.  A closure is
acceptable only if it specifies how \(\mathcal C_j\) is estimated and if the
result is checked against the direct root-solve integration.

\section{Derivative Observables}

Differentiate the charge-balance equation only after the state and root solve
are defined.  For the isothermal case,
\begin{equation}
\frac{\dd V_n}{\dd q}
=
\frac{
Q_{\cell}
-\sum_{j=1}^{N_p}Q_{j,\tot}\frac{\dd \xi_j}{\dd q}
}{
\frac{\partial Q_{\bg}}{\partial V}(V_n,T)
}.
\label{eq:dVdq_iso_draft}
\end{equation}
With temperature variation, include the \(\partial Q_{\bg}/\partial T\) term
and any declared temperature dependence of the fitted parameters.

The apparent-voltage derivative is
\begin{equation}
\frac{\dd V_{n,\app}}{\dd q}
=
\frac{\dd V_n}{\dd q}
+s_I\frac{\dd}{\dd q}\bigl(|I|R_n(q,T,|I|)\bigr).
\label{eq:dVappdq_draft}
\end{equation}

The ICA and DVA observables are based on external charge:
\begin{align}
\frac{\dd Q_{\ext}}{\dd V_{n,\app}}
&=
\frac{Q_{\cell}}{\dd V_{n,\app}/\dd q},
\label{eq:ica_draft}\\
\frac{\dd V_{n,\app}}{\dd Q_{\ext}}
&=
\frac{1}{Q_{\cell}}
\frac{\dd V_{n,\app}}{\dd q}.
\label{eq:dva_draft}
\end{align}

\section{Validation Gates}

The model is admissible only when the following checks are satisfied:
\begin{enumerate}[label=\arabic*.]
\item The normalized charge-balance residual
  \(\mathcal G(V_n;q,T,\bm{\xi};\theta)/Q_{\cell}\) is small at all evaluated
  points.
\item The charge-balance root remains inside the allowed voltage interval.
\item The slope floor in Eq.~\eqref{eq:bg_slope_floor_draft} is not violated.
\item The transition states remain in their declared physical domain.
\item The capacity closure is consistent at the endpoints of the analysis
  window.
\item The denominator in Eq.~\eqref{eq:ica_draft} does not cross an
  unphysical singularity.
\item Any reference-closure approximation is compared against the direct
  root-solve integration.
\end{enumerate}

\section{Interfaces to Later Chapters}

Chapter 2 may use \(T\), \(V_n\), \(V_{n,\app}\), \(\xi_j\), and
\(\dd\xi_j/\dd t\) for heat coupling.  Chapter 3 may use \(V_{n,\drive}\),
\(\mathcal A_j\), and \(k_j\) for electrochemical kinetic extensions.  Chapter
4 may assemble residuals, constraints, and observables into a fitting system.
Chapter 5 may introduce hysteresis or structural memory, but it must preserve
the charge-balance residual and the distinction between \(V_n\), \(V_{n,\app}\),
and \(V_{n,\drive}\).

\begin{thebibliography}{9}

\bibitem{Lee2017Field}
K. Lee, S. Lee, C. H. Choi, and S. Lee,
``Effects of external electric field and anisotropic long-range reactivity on
charge separation probability,''
\textit{J. Chem. Phys.} \textbf{147}, 144111 (2017).
doi:10.1063/1.5000882.

\bibitem{Lee2011Propagator}
S. Lee, C. Y. Son, J. Sung, and S.-H. Chong,
``Communication: Propagator for diffusive dynamics of an interacting molecular
pair,''
\textit{J. Chem. Phys.} \textbf{134}, 121102 (2011).
doi:10.1063/1.3565476.

\bibitem{Son2013Rates}
C. Y. Son, J. Kim, J.-H. Kim, J. S. Kim, and S. Lee,
``An accurate expression for the rates of diffusion-influenced bimolecular
reactions with long-range reactivity,''
\textit{J. Chem. Phys.} \textbf{138}, 164123 (2013).
doi:10.1063/1.4802584.

\end{thebibliography}

\end{document}
